\documentclass[11pt]{article}
\usepackage[margin=1.1in]{geometry}
\usepackage{amsmath,amssymb,amsthm}
\usepackage{booktabs}
\usepackage{array}
\usepackage{graphicx}
\usepackage[hidelinks]{hyperref}
\usepackage{xcolor}
\usepackage{parskip}

\graphicspath{{../reports/heatmaps/gen300_20260722/}{../reports/heatmaps/pairfloat_20260723/}}
\newcommand{\code}[1]{\texttt{#1}}
\newcommand{\gf}{\mathrm{GF}(2)}
\newtheorem{theorem}{Theorem}

\title{\textbf{Reducing the Progress-Diagonal Ridge}\\[2pt]
\large Coverage experiments and the use-point wall}
\author{local\_mixing project}
\date{July 23, 2026 \\ \small branch \code{ssg-gen-mix-clean} --- tools
\code{fmix --gen-snap-every} (\code{49cee4c8}) and
\code{experimental/mixing/pairfloat\_exp.rs} (\code{5a9f449c})}

\begin{document}
\maketitle

\begin{abstract}
\noindent Local mixing rewrites a small circuit $C$ into a large equivalent
$G$ that should hide $C$'s computation. The sharpest leak is the
\emph{progress diagonal}: a bright ridge in the affine-reconstruction heatmap
between $C$ and $G$, showing how far through the original computation any
prefix of $G$ has reached. We ask whether the ridge depth can be driven to
$0.05$ or below. After clearing two non-levers (generation dose and affine
moves), we test a family of \emph{free nonlinear brackets} --- inserted
identity pairs floated apart --- across four coverage strategies of increasing
sophistication. All fail, the most careful one (gap-tiling) failing
\emph{backwards}: more masking sharpens the ridge. The failures are unanimous
and their cause is a single structural fact, the \emph{use-point theorem}: the
ridge lives where operands are consumed, and no relocation-based masking
touches that location. The only remaining route is computation on encoded
operands (share-native / mask-aware gates).
\end{abstract}

\section{The question}

The reconstruction heatmap $H(i,j)$ is the best-effort $\gf$-affine
reconstruction error of $C$'s state after $i$ gates from \emph{all} wires of
$G$'s state after $j$ gates. It shows a bright ridge along the line where $G$'s
computational progress matches $C$'s; the ridge reveals the temporal location
of the computation. We read the heatmap by the \emph{ridge measure}
(\code{reports/plot\_hmap\_ridge.py}): \code{depth} (ridge prominence,
$0=$ flat/ideal, ${\sim}0.5=$ a maximally bright ridge), \code{rho} (Spearman
of $C$-prefix vs.\ ridge column; $\code{rho}=1$ is a clean monotone diagonal),
a permutation $z$-score, and \code{contrast}. The ridge has survived every
mixing method tried, always at $\code{rho}=1$, with \code{depth} floored near
$0.23$--$0.35$. This note asks whether \code{depth} can reach $0.05$.

\section{Two non-levers}

\paragraph{Generation dose.} \code{fmix} phase A drives each gate through a
target number of database re-encodings (``generation''). We ran a fresh
gadgetized $n=128$ sliced sandwich (85{,}058 gates) to generation $300$ ---
six times the benchmarked dose --- snapshotting at each generation multiple of
$50$ via the new \code{--gen-snap-every} flag (Fig.~\ref{fig:gen300}). The
ridge only sharpens toward a floor:

\begin{center}
\begin{tabular}{lccccccc}
\toprule
generation & 0 & 50 & 100 & 150 & 200 & 250 & 300\\
\midrule
gates & 85{,}058 & 179{,}437 & 218{,}809 & 265{,}096 & 317{,}289 & 377{,}912 & 450{,}318\\
ridge depth & 0.348 & 0.273 & 0.261 & 0.254 & 0.246 & 0.240 & 0.236\\
rho & 1.00 & 1.00 & 1.00 & 1.00 & 1.00 & 1.00 & 1.00\\
\bottomrule
\end{tabular}
\end{center}

Decrements shrink geometrically toward ${\sim}0.22$--$0.23$. Six times the dose
past generation~50 bought a $14\%$ depth reduction at $2.5\times$ the size.

\paragraph{Affine moves.} The measure is invariant to the affine part of any
encoding (CNOT, negation, XOR-sharing). Every conjugation twist in \code{fmix}
is affine, so it \emph{provably} cannot move the degree-1 ridge, at any dose.
What is left is genuinely nonlinear re-encoding, which database replacements
supply only inside ${\le}5$-gate windows; the $0.23$ floor is the saturation of
that thin nonlinear layer.

\begin{figure}[t]
\centering
\includegraphics[width=\linewidth]{gen300_ridge_progression.png}
\caption{Generation-dose progression, gen $0\to300$. The diagonal stays at
$\code{rho}=1$; only the depth shallows, toward a floor near $0.23$.}
\label{fig:gen300}
\end{figure}

\section{Free nonlinear brackets}

Insert an identical pair $g\cdot g$ (the identity) into $C$, then float the two
copies apart by adjacent swaps through gates they commute with. Correctness is
exact and preserved bit-for-bit. Between the separated copies exactly one $g$
is applied, so every prefix cut there reads the $g$-image of the true state ---
the target of $g$ appears at degree $1+(\#\text{controls})$. A separated pair
is a \emph{free nonlinear bracket}: no rewriting, only relocation. The tool
\code{pairfloat\_exp} builds a fixed $C$ (3000 g57 gates, 128 wires) and tests
families of this move, verifying equivalence on 200 random inputs per output.

\section{Four attempts}

\paragraph{4.1 Insertion count $k$ (g57 pairs, commuting float).}
\begin{center}
\begin{tabular}{lcccccccccc}
\toprule
$k$ & 0 & 1 & 2 & 3 & 4 & 5 & 8 & 11 & 16 & 24\\
depth & 0.474 & 0.395 & 0.377 & 0.370 & 0.367 & 0.364 & 0.359 & 0.357 & 0.357 & 0.351\\
\bottomrule
\end{tabular}
\end{center}
The first pair helps; the curve then flattens onto ${\sim}0.35$. \emph{Self-dilution}:
straddling brackets per cut $=\text{pairs}\times\text{span}/\text{size}
=84000/(1+2k)$, saturating at $\text{span}/2\approx42$ ($\sim27\%$ of wires) as
$k\to\infty$ --- each pair inflates the circuit it must cover. A model
$\code{depth}\approx\code{depth}_0\,(1-\text{coverage})$ fits within ${\sim}0.01$.

\paragraph{4.2 Cross budget $B$ (float past colliders by conjugation).} Free
floating stops after ${\sim}42$ slots. Crossing up to $B$ colliders via the
exact splitting rules (\code{engine::rules} R1/R2/R3) grows span, size, and
loses half the floaters as colliders:
\begin{center}
\begin{tabular}{lccccc}
\toprule
$B$ (at $k=3$) & 0 & 1 & 2 & 4 & 8\\
\midrule
depth & 0.369 & 0.343 & 0.333 & 0.330 & 0.329\\
mean travel & 41 & 120 & 192 & 302 & 405\\
gates & 21k & 53k & 77k & 105k & 129k\\
\bottomrule
\end{tabular}
\end{center}
Span rose $10\times$; depth saturated at $0.33$; the whole gain is banked by
$B=1$. Crossing is where the free bracket becomes priced, and the price outruns
the benefit at once.

\paragraph{4.3 Adaptive wire selection.} Choosing each pair's wires for
maximal mobility (quiet target, unwritten controls) raised span $7.5$--$50\times$
and made depth \emph{worse}:
\begin{center}
\begin{tabular}{lccccc}
\toprule
adaptive & $k{=}1$ & $k{=}3$ & $k{=}5$ & $+B{=}4$ & $+B{=}8$\\
\midrule
depth & 0.437 & 0.427 & 0.423 & 0.402 & 0.396\\
mean travel & 150 & 310 & 432 & 1486 & 2039\\
\bottomrule
\end{tabular}
\end{center}
\emph{Target concentration}: a wire is mobile because nothing near it is
active, so it stays quiet for hundreds of cuts and every insertion samples the
same tiny pool. Distinct masked wires collapse to ${\sim}10$--$20$ of $128$,
dropping coverage below random's. Mobility and diversity are in tension.

\paragraph{4.4 Gap-tiling --- the decisive result.} Invert the assignment: for
each wire $w$, every maximal stretch of $C$ between consecutive touches of $w$
is a bracket slot; insert $s$ distinct pairs targeting $w$ there (controls
untargeted in the gap, so the pair commutes with the interior), floated to the
gap ends. This forces diversity over all wires and masks every wire at every
cut \emph{except its own use points}. Material cost matched prediction
($\times6.9/12.9/24.8$ for $s=1/2/4$); the predicted depth was $0.25/0.12/0.05$.
The result was the opposite (Fig.~\ref{fig:gaptile}):
\begin{center}
\begin{tabular}{lccc}
\toprule
gap-tiling $s$ & 1 & 2 & 4\\
\midrule
depth & 0.364 & 0.375 & 0.398\\
contrast & 3.26 & 3.26 & 3.26\\
\bottomrule
\end{tabular}
\end{center}
Stacking masks makes the ridge \emph{sharper}. The plate shows why: the
off-diagonal background hides more uniformly while the diagonal is untouched
and stands out more by contrast.

\begin{figure}[t]
\centering
\includegraphics[width=\linewidth]{pairfloat_gaptile.png}
\caption{Gap-tiling $s=1,2,4$ (rightmost three) against $C$-vs-$C$ and the
random-insertion baseline. Depth \emph{rises} with stacking: masking between
uses cannot reach the diagonal.}
\label{fig:gaptile}
\end{figure}

\section{Conclusion: the use-point wall}

\begin{theorem}[use-point, empirical]
Between-use bracket masking --- in any form (random $k$, cross budget,
adaptive mobility, gap-tiling) --- cannot reduce the degree-1 progress ridge
below ${\sim}0.33$. It targets the wrong location.
\end{theorem}

\noindent Gap-tiling masks each wire only strictly between its uses (floaters
stop at the touch gates), but the ridge at $C$-row $i$ aligns with the
$G$-column where gate $i$'s operands are \emph{used}, and those are left
exposed. The diagonal \emph{is} the operand-use locus, so between-use masking
cannot reach it; hiding the idle background only sharpens the aligned column by
contrast. The small reductions the earlier modes obtained ($0.47\to0.33$) came
precisely from brackets that \emph{accidentally straddled} use points while
floating blindly --- gap-tiling removes those straddles and does \emph{worse}
than random.

This is the same wall the deferred-mask work
(\code{docs/NONLINEAR\_SHARE\_ENCODING}) hit from the other side, where a peek
un-masked each operand at its use point and re-exposed the diagonal. Two
independent attacks converge: the ridge cannot be moved by relocation.

\section{What remains}

To reduce the ridge one must mask a wire \emph{at the instant it is consumed}:
compute a gate on the encoded operand $w\oplus\text{mask}$ without ever
reconstructing $w$. That is share-native / mask-aware computation --- what the
legacy g57 gadget does (it computes on shares and never reconstructs an
operand, which is why it hides at degree~2). The open construction is a
mask-aware CG that folds the mask correction into the gate; its cost is the
answer-extraction interface, not gate count. Coverage-only approaches, and
their ${\sim}\times25$-material price for full tiling, are a closed dead end.

\end{document}
